%% ============================================================
%% ENTREGABLE — TE3001B Fundamentación de Robótica
%% Teleoperación con Sensado de Fuerza: xArm Lite 6 (Maestro-Esclavo)
%% Formato IEEE Transactions — Plantilla para Alumnos
%% Prof. Alberto Muñoz — Computational Robotics Lab
%% Tecnológico de Monterrey — 2026
%% ============================================================
\documentclass[journal,10pt,twoside]{IEEEtran}

%% ── Paquetes base ────────────────────────────────────────────
\usepackage[utf8]{inputenc}
\usepackage[T1]{fontenc}
\usepackage[spanish,es-nodecimaldot]{babel}
\usepackage{amsmath,amssymb,bm}
\usepackage{graphicx}
\usepackage{xcolor}
\usepackage{booktabs}
\usepackage{multirow}
\usepackage{array}
\usepackage{colortbl}
\usepackage{caption}
\usepackage{subcaption}
\usepackage{listings}
\usepackage{tcolorbox}
\usepackage{enumitem}
\usepackage{tikz}
\usepackage{pgfplots}
\usepackage{url}
\usepackage{hyperref}   %% hyperref DESPUÉS de xcolor para evitar conflictos
\pgfplotsset{compat=1.18}
\usetikzlibrary{arrows.meta,shapes,positioning,calc,backgrounds}
\tcbuselibrary{breakable,skins}

%% ── Colores institucionales ──────────────────────────────────
%% IMPORTANTE: definir colores ANTES de \hypersetup
\definecolor{tec_blue}{RGB}{0,68,124}
\definecolor{tec_lightblue}{RGB}{0,140,200}
\definecolor{tec_red}{RGB}{190,30,45}
\definecolor{tec_gray}{RGB}{100,100,100}
\definecolor{code_bg}{RGB}{245,247,250}
\definecolor{code_kw}{RGB}{0,100,180}
\definecolor{code_str}{RGB}{180,50,20}
\definecolor{code_cmt}{RGB}{80,140,80}
\definecolor{box_todo}{RGB}{255,243,205}
\definecolor{box_todo_border}{RGB}{200,140,0}
\definecolor{box_rubric}{RGB}{232,245,233}
\definecolor{box_rubric_border}{RGB}{67,160,71}
\definecolor{box_theory}{RGB}{227,242,253}
\definecolor{box_theory_border}{RGB}{25,118,210}

%% hypersetup aquí, DESPUÉS de definir tec_blue ───────────────
\hypersetup{colorlinks=true, linkcolor=tec_blue, urlcolor=tec_blue,
            citecolor=tec_blue}

%% ── Cajas de tcolorbox ───────────────────────────────────────
\newtcolorbox{todobox}[1][Completar]{
  colback=box_todo, colframe=box_todo_border,
  title={\textbf{$\triangleright$ #1}},
  fonttitle=\bfseries\color{tec_gray},
  breakable, enhanced, left=4pt, right=4pt
}
\newtcolorbox{rubricbox}[1][Rúbrica]{
  colback=box_rubric, colframe=box_rubric_border,
  title={\textbf{\checkmark\ Criterio de Evaluación — #1}},
  fonttitle=\bfseries\small\color{box_rubric_border},
  breakable, enhanced, left=4pt, right=4pt
}
\newtcolorbox{theorybox}[1][Marco Teórico]{
  colback=box_theory, colframe=box_theory_border,
  title={\textbf{(i) Contexto: #1}},
  fonttitle=\bfseries\small\color{tec_blue},
  breakable, enhanced, left=4pt, right=4pt
}

%% ── Estilo de código ─────────────────────────────────────────
\lstdefinestyle{python}{
  language=Python,
  backgroundcolor=\color{code_bg},
  basicstyle=\ttfamily\scriptsize,
  keywordstyle=\color{code_kw}\bfseries,
  stringstyle=\color{code_str},
  commentstyle=\color{code_cmt}\itshape,
  numbers=left, numberstyle=\tiny\color{tec_gray},
  numbersep=6pt, frame=single, tabsize=4,
  showstringspaces=false, breaklines=true,
  morekeywords={import,from,as,def,class,return,True,False,None,self},
  literate={á}{{\'{a}}}1 {é}{{\'{e}}}1 {í}{{\'{i}}}1
           {ó}{{\'{o}}}1 {ú}{{\'{u}}}1 {ñ}{{\~{n}}}1,
}
\lstset{style=python}

%% ── Macros de ayuda ──────────────────────────────────────────
\newcommand{\answerline}[1][3cm]{\underline{\hspace{#1}}}
\newcommand{\pts}[1]{\textcolor{tec_red}{\textbf{[#1 pts]}}}
\newcommand{\TODO}[1]{\begin{todobox}[Respuesta del equipo]
  \textit{#1}
\end{todobox}}
\newcommand{\FIGURE}[1]{\begin{todobox}[Insertar figura]
  \centering
  %% \resizebox escala el tikzpicture al ancho disponible
  \resizebox{\linewidth}{!}{%
    \begin{tikzpicture}
      \draw[tec_gray, dashed, line width=1pt] (0,0) rectangle (8,4);
      \node[tec_gray, font=\small] at (4,2) {#1};
    \end{tikzpicture}%
  }
\end{todobox}}

%% ── Metadatos del artículo ───────────────────────────────────
\begin{document}

\title{Teleoperación con Sensado de Fuerza Física\\
  entre dos xArm Lite 6:\\
  Diseño, Implementación y Evaluación Experimental}

\author{%
  \IEEEauthorblockN{%
    Nombre Alumno 1\IEEEauthorrefmark{1},\;
    Nombre Alumno 2\IEEEauthorrefmark{1},\;
    Nombre Alumno 3\IEEEauthorrefmark{1}%
  }\\
  \IEEEauthorblockA{%
    \IEEEauthorrefmark{1}%
    Computational Robotics Lab, Tecnológico de Monterrey,
    Monterrey, N.L., México\\
    \texttt{\{matrícula1, matrícula2, matrícula3\}@tec.mx}
  }
}

\maketitle

%% ── Abstract ─────────────────────────────────────────────────
\begin{abstract}
%% tcolorbox NO puede usarse dentro del abstract de IEEEtran.
%% Reemplazar este párrafo con el resumen real (150--200 palabras).
%% Estructura sugerida: (1) problema, (2) arquitectura maestro-esclavo
%% con xArm Lite 6, (3) estrategia de sensado de fuerza, (4) resultados
%% experimentales (error de posición, fuerza de contacto, latencia),
%% (5) conclusiones principales. No incluyan fórmulas en el abstract.
\textit{[Completar: escriban aquí el resumen de 150--200 palabras.]}
\end{abstract}

\begin{IEEEkeywords}
  teleoperación, sensado de fuerza, xArm Lite 6, control de impedancia,
  retroalimentación háptica, \answerline[2cm], \answerline[2cm]
\end{IEEEkeywords}

%% ============================================================
\section{Introducción}
%% ============================================================

\IEEEPARstart{L}{a} teleoperación con retroalimentación de fuerza
constituye un paradigma fundamental en robótica cuando se requiere que
un operador humano interactúe con entornos remotos o peligrosos sin
presencia física~\cite{murray1994}. El xArm Lite 6, desarrollado por
UFACTORY~\cite{xarm_github}, es un manipulador de 6 grados de libertad
(DOF) con capacidad de control de torque que lo hace idóneo para experimentos
de sensado de fuerza por torque articular.

\begin{theorybox}[Motivación]
  El reto central de este proyecto es cerrar el lazo de fuerza entre
  dos brazos robóticos: el robot \textbf{Maestro}, operado por el humano,
  y el robot \textbf{Esclavo}, que interactúa con el entorno.
  La fuerza de contacto estimada en el esclavo debe reflejarse como
  torque en las articulaciones del maestro, creando la sensación
  háptica que guía al operador.
\end{theorybox}

\subsection{Contexto y Motivación}
\TODO{Describan en 2--3 párrafos por qué es relevante la teleoperación
con sensado de fuerza. Mencionen al menos DOS aplicaciones industriales
o médicas concretas donde se utiliza esta tecnología. Citen al menos
3 referencias del estado del arte (artículos IEEE, libros de texto).}

\subsection{Contribuciones del Trabajo}
\TODO{Enumeren de forma concisa (lista con viñetas) las 3--5 contribuciones
específicas de su implementación. Por ejemplo: diseño del protocolo de
comunicación, estrategia de estimación de fuerza elegida, experimentos
realizados, etc.}

\subsection{Estructura del Artículo}
El resto de este artículo se organiza como sigue:
la Sección~\ref{sec:background} presenta los fundamentos teóricos;
la Sección~\ref{sec:architecture} describe la arquitectura del sistema;
la Sección~\ref{sec:force} detalla la estrategia de sensado de fuerza;
la Sección~\ref{sec:control} expone el esquema de control;
la Sección~\ref{sec:implementation} documenta la implementación con xArm SDK;
la Sección~\ref{sec:experiments} reporta los resultados experimentales;
la Sección~\ref{sec:discussion} discute los hallazgos;
y la Sección~\ref{sec:conclusions} concluye el trabajo.

%% ============================================================
\section{Marco Teórico}
\label{sec:background}
%% ============================================================

\subsection{Cinemática y Jacobiano del xArm Lite 6}

El xArm Lite 6 es un manipulador de $n=6$ DOF con articulaciones
rotacionales. Su cinemática directa se describe mediante la convención
Denavit-Hartenberg (DH) o Product of Exponentials (PoE).

\begin{theorybox}[Jacobiano Geométrico]
  El Jacobiano $\bm{J}(\bm{q}) \in \mathbb{R}^{6\times6}$ mapea
  velocidades articulares a velocidades cartesianas del efector final:
  \[
    \begin{pmatrix} \bm{v} \\ \bm{\omega} \end{pmatrix} = \bm{J}(\bm{q})\,\dot{\bm{q}}
  \]
  Su transpuesta mapea fuerzas/torques cartesianos a torques articulares:
  \[
    \bm{\tau} = \bm{J}^\top(\bm{q})\,\bm{F}_{ext}
  \]
  Esta relación es la base del \emph{torque-based force sensing}.
\end{theorybox}

\subsubsection*{Pregunta~T1 \pts{8}}
\begin{enumerate}[label=(\alph*)]
  \item Expliquen la diferencia entre el \textbf{Jacobiano analítico}
        y el \textbf{Jacobiano geométrico} para el xArm Lite 6.
        ¿Cuál utilizan en su implementación y por qué?
  \item Identifiquen y describan \textbf{dos configuraciones singulares}
        del xArm Lite 6 a partir de su geometría DH. ¿Qué ocurre con la
        estimación de fuerza $\hat{\bm{F}} = (\bm{J}^\top)^+\bm{\tau}$
        cerca de esas singularidades?
  \item ¿Cómo detectan programáticamente que el robot está cerca de
        una singularidad? Escriban la condición en términos del
        número de condición $\kappa(\bm{J})$.
\end{enumerate}
\TODO{Respondan aquí la Pregunta T1. Incluyan las expresiones
matemáticas correspondientes y una figura con las configuraciones
singulares identificadas en el xArm Lite 6.}

\FIGURE{Fig. T1: Configuraciones singulares del xArm Lite 6.
Señalen las articulaciones involucradas.}

\subsection{Dinámica del Manipulador}

La dinámica de un robot rígido de $n$ DOF se expresa como:
\begin{equation}
  \bm{M}(\bm{q})\ddot{\bm{q}} + \bm{C}(\bm{q},\dot{\bm{q}})\dot{\bm{q}}
  + \bm{g}(\bm{q}) = \bm{\tau} - \bm{J}^\top\bm{F}_{ext}
  \label{eq:dynamics}
\end{equation}
donde $\bm{M}$ es la matriz de inercia, $\bm{C}$ la matriz de
Coriolis/centrífuga y $\bm{g}$ el vector gravitacional.

\subsubsection*{Pregunta~T2 \pts{10}}
\begin{enumerate}[label=(\alph*)]
  \item Partiendo de la Ec.~\eqref{eq:dynamics}, \textbf{deriven} la
        ley de control de par calculado (Computed Torque Control, CTC)
        y demuestren que produce la dinámica de error lineal:
        \[
          \ddot{\bm{e}} + K_v\dot{\bm{e}} + K_p\bm{e} = \bm{0}
        \]
        ¿Qué condiciones sobre $K_p$ y $K_v$ garantizan
        estabilidad asintótica global?
  \item Para el xArm Lite 6 operando en modo \texttt{SERVO} de
        la SDK, el bucle de control interno trabaja a \answerline[1.5cm]~Hz.
        Expliquen cómo esto limita la elección de $K_p$ y $K_v$.
  \item Expresen los polos del sistema cerrado en términos de
        $\omega_n$ (frecuencia natural) y $\zeta$ (amortiguamiento).
        ¿Qué valor de $\zeta$ usaron y por qué?
\end{enumerate}
\TODO{Respondan aquí la Pregunta T2. Incluyan la derivación paso a paso
y una gráfica de la respuesta al escalón del controlador sintonizado
con sus valores de $K_p$, $K_v$.}

\FIGURE{Fig. T2: Respuesta al escalón del controlador CTC con sus
ganancias $K_p$ y $K_v$. Señalen $\omega_n$, $\zeta$, tiempo de
establecimiento y sobreimpulso.}

\subsection{Control de Impedancia}

El control de impedancia~\cite{hogan1985} modela el efector final como
un sistema masa-resorte-amortiguador virtual:
\begin{equation}
  \bm{F}_{ctrl} = \bm{K}_d(\bm{x}_d - \bm{x})
                + \bm{B}_d(\dot{\bm{x}}_d - \dot{\bm{x}})
                + \bm{M}_d(\ddot{\bm{x}}_d - \ddot{\bm{x}})
  \label{eq:impedance}
\end{equation}

\subsubsection*{Pregunta~T3 \pts{8}}
\begin{enumerate}[label=(\alph*)]
  \item Para la Ec.~\eqref{eq:impedance}, analicen el efecto de
        aumentar $\bm{K}_d$ en: (i) la sensación háptica del operador,
        (ii) la estabilidad con retardo $T_d$, (iii) la velocidad
        de respuesta ante contacto.
  \item Dado un retardo de red $T_d = 20$~ms medido en su laboratorio,
        calculen la rigidez virtual máxima $K_{d,max}$ que garantiza
        estabilidad usando el criterio de pasividad de Colgate~\cite{colgate1993}:
        \[
          K_{d,max} < \frac{\pi}{2\,T_d\,\omega_c}
        \]
        donde $\omega_c$ es el ancho de banda del controlador interno
        del xArm. Justifiquen el valor de $\omega_c$ que usan.
  \item ¿Usaron impedancia en espacio articular o cartesiano?
        Justifiquen su elección para la tarea de teleoperación.
\end{enumerate}
\TODO{Respondan aquí la Pregunta T3. Incluyan los cálculos numéricos
y una tabla comparativa de comportamiento con distintos valores de
$K_d$ probados.}

%% ============================================================
\section{Arquitectura del Sistema}
\label{sec:architecture}
%% ============================================================

\subsection{Visión General}

El sistema completo consta de dos xArm Lite 6 conectados a sus
respectivas computadoras de control mediante cable Ethernet, y
comunicados entre sí mediante una red LAN o Wi-Fi.

\subsubsection*{Pregunta~A1 \pts{10}}
\begin{enumerate}[label=(\alph*)]
  \item Elaboren un \textbf{diagrama de bloques} completo del sistema
        maestro-esclavo que incluya: xArm Maestro, PC Maestro,
        protocolo de comunicación, PC Esclavo, xArm Esclavo,
        sensor/estimador de fuerza, y lazo de retroalimentación háptica.
        Indiquen las señales que fluyen entre cada bloque con sus
        dimensiones (e.g., $\bm{q}\in\mathbb{R}^6$, $\bm{F}\in\mathbb{R}^3$).
  \item Describan el \textbf{flujo de datos} en un ciclo de control:
        ¿qué señales se leen del robot Maestro? ¿qué se envía al Esclavo?
        ¿qué se retroalimenta al Maestro?
  \item ¿Cuál es la frecuencia de actualización de su lazo de control?
        Justifiquen por qué eligieron esa frecuencia en relación con
        las capacidades del xArm SDK (\texttt{set\_servo\_angle} vs.
        \texttt{set\_position}).
\end{enumerate}

\FIGURE{Fig. A1: Diagrama de bloques del sistema de teleoperación
maestro-esclavo con xArm Lite 6. Señalen todas las señales con sus
dimensiones.}

\TODO{Respondan aquí la Pregunta A1.}

\subsection{Protocolo de Comunicación}

\subsubsection*{Pregunta~A2 \pts{8}}
\begin{enumerate}[label=(\alph*)]
  \item Comparen el uso de \textbf{UDP vs TCP} para la comunicación
        en tiempo real entre los dos PCs. ¿Cuál eligieron y por qué?
        Completen la siguiente tabla con los resultados medidos:

        \smallskip
        \resizebox{\linewidth}{!}{%
        \begin{tabular}{lcc}
          \toprule
          \textbf{Métrica} & \textbf{UDP} & \textbf{TCP} \\
          \midrule
          RTT promedio (ms)      & \answerline[2cm] & \answerline[2cm] \\
          Jitter (ms)            & \answerline[2cm] & \answerline[2cm] \\
          Paquetes perdidos (\%) & \answerline[2cm] & N/A \\
          Throughput (KB/s)      & \answerline[2cm] & \answerline[2cm] \\
          \bottomrule
        \end{tabular}}
        \smallskip

  \item Definan el formato del \textbf{paquete de datos} que intercambian.
        Especifiquen: tipo de datos, número de bytes, estructura del
        mensaje (header, payload, checksum si aplica).
  \item Con un RTT medido de $T_{RTT}$, ¿cuál es el retardo efectivo
        $T_d = T_{RTT}/2$ en su sistema? ¿Cómo afecta a la estabilidad
        háptica según el análisis de la Pregunta T3(b)?
\end{enumerate}

\TODO{Respondan aquí la Pregunta A2. Incluyan la captura de pantalla
de \texttt{ping} y de su herramienta de test de red entre las dos PCs.}

\FIGURE{Fig. A2: Histograma de RTT medido entre PC Maestro y
PC Esclavo. Incluyan estadísticas: media, desviación estándar, percentil 99.}

%% ============================================================
\section{Estrategia de Sensado de Fuerza}
\label{sec:force}
%% ============================================================

\begin{theorybox}[¿Cómo «siente» la fuerza un xArm sin sensor F/T externo?]
  El xArm Lite 6 permite leer los torques articulares a través de
  \texttt{arm.get\_joint\_states()} o mediante el modo de control de
  torque. A partir de los torques medidos $\bm{\tau}_{meas}$ y el
  modelo dinámico del robot, es posible \emph{estimar} la fuerza
  externa en el efector final sin instalar un sensor F/T dedicado.
  Esta técnica se denomina \textbf{sensado de fuerza basado en torque articular}.
\end{theorybox}

\subsection{Estimación de Fuerza por Torque Articular}

La fuerza externa en el efector se estima como:
\begin{equation}
  \hat{\bm{F}}_{ext} = \left(\bm{J}^\top\right)^+
    \underbrace{\left(\bm{\tau}_{meas} - \bm{\tau}_{model}\right)}_{\Delta\bm{\tau}}
  \label{eq:force_estimate}
\end{equation}
donde $\bm{\tau}_{model} = \bm{M}(\bm{q})\ddot{\bm{q}} +
\bm{C}\dot{\bm{q}} + \bm{g}(\bm{q})$ es el torque predicho por el modelo.

\subsubsection*{Pregunta~F1 \pts{12}}
\begin{enumerate}[label=(\alph*)]
  \item Expliquen \textbf{tres fuentes de error} en la estimación de
        $\hat{\bm{F}}_{ext}$ usando la Ec.~\eqref{eq:force_estimate}
        para el xArm Lite 6. Proporcionen una estimación cuantitativa
        del error que introduce cada una.
  \item En su implementación, ¿compensaron la gravedad antes de estimar
        la fuerza? Escriban la expresión de $\bm{g}(\bm{q})$ utilizada
        o describan cómo la obtuvieron del modelo URDF del xArm.
  \item Apliquen un \textbf{filtro pasa-bajas} a $\hat{\bm{F}}_{ext}$.
        Escriban la función de transferencia discreta del filtro
        utilizado, justifiquen la frecuencia de corte elegida
        y muestren su respuesta en frecuencia.
  \item ¿Cuál es el umbral de fuerza $F_{th}$ que eligieron para
        detectar contacto? Justifiquen en términos de la relación
        señal-ruido (SNR) medida con el robot en reposo.
\end{enumerate}

\TODO{Respondan aquí la Pregunta F1. Incluyan la expresión explícita
de cada término de la Ec.~\eqref{eq:force_estimate} para su
configuración específica de laboratorio.}

\FIGURE{Fig. F1: (a) Torques articulares medidos vs. predichos por
el modelo en una trayectoria de prueba. (b) Fuerza estimada
$\hat{\bm{F}}_{ext}(t)$ sin y con filtro pasa-bajas.}

\subsection{Alternativas de Sensado Consideradas}

\subsubsection*{Pregunta~F2 \pts{8}}
\begin{enumerate}[label=(\alph*)]
  \item Comparen las siguientes estrategias de sensado de fuerza
        para el xArm Lite 6. Completen la tabla:

\smallskip
\resizebox{\linewidth}{!}{%
\begin{tabular}{p{2.4cm}p{2cm}p{2cm}p{1cm}}
  \toprule
  \textbf{Estrategia} & \textbf{Ventajas} & \textbf{Desventajas} & \textbf{Costo} \\
  \midrule
  Torque articular (sin sensor externo)
    & \answerline[1.8cm] & \answerline[1.8cm] & Bajo \\
  Sensor F/T en muñeca (DIY)
    & \answerline[1.8cm] & \answerline[1.8cm] & Alto \\
  Visión por computadora (force from deformation)
    & \answerline[1.8cm] & \answerline[1.8cm] & Medio \\
  Galgas extensométricas en el efector
    & \answerline[1.8cm] & \answerline[1.8cm] & Medio \\
  \bottomrule
\end{tabular}}
\smallskip

  \item Justifiquen la estrategia que implementaron indicando
        la razón principal (costo, disponibilidad de hardware,
        precisión requerida, tiempo de implementación).
  \item Si tuvieran un sensor en la muñeca del
        esclavo, ¿cómo modificaron su pipeline de fuerza?
        Describan los cambios en software y hardware.
\end{enumerate}

\TODO{Respondan aquí la Pregunta F2.}

%% ============================================================
\section{Esquema de Control}
\label{sec:control}
%% ============================================================

\subsection{Control del Robot Esclavo}

\subsubsection*{Pregunta~C1 \pts{10}}
\begin{enumerate}[label=(\alph*)]
  \item Describan la ley de control implementada en el robot Esclavo.
        Escriban la expresión matemática completa e indiquen qué modo
        de control del xArm SDK utilizaron
        (\texttt{set\_servo\_angle}, \texttt{set\_position},
        o \texttt{set\_servo\_cartesian}).
        ¿Por qué eligieron ese modo?
  \item Muestren el diagrama de flujo (flowchart) del ciclo de control
        del esclavo: lectura de posición del maestro $\rightarrow$
        cálculo de referencia $\rightarrow$ envío al xArm $\rightarrow$
        lectura de torques $\rightarrow$ estimación de fuerza
        $\rightarrow$ envío de feedback.
  \item Introdujeron algún escalado de posición/velocidad entre
        maestro y esclavo (factor $\alpha$)?
        \[
          \bm{x}_{slave}^{ref} = \alpha\,\bm{x}_{master}
        \]
        Justifiquen el valor de $\alpha$ elegido considerando
        el espacio de trabajo de ambos robots.
  \item Implementaron alguna \textbf{limitación de fuerza}
        (\textit{force clamping}) para proteger el robot esclavo?
        Describan la condición de saturación.
\end{enumerate}

\TODO{Respondan aquí la Pregunta C1. Incluyan el fragmento de código
Python que implementa el ciclo de control del esclavo.}

\begin{lstlisting}[caption={Ciclo de control del robot Esclavo — Completar con su implementación real},
                   label={lst:slave}]
# ── Completar con el código real de su implementación ──
# Ciclo principal del robot Esclavo
import time
from xarm.wrapper import XArmAPI

arm_slave = XArmAPI('192.168.1.XXX')  # <-- IP real del esclavo
arm_slave.motion_enable(enable=True)
arm_slave.set_mode(?)  # <-- modo elegido y justificación
arm_slave.set_state(0)

while running:
    # 1. Recibir posición/referencia del Maestro vía UDP
    ref = receive_from_master()   # <-- implementar

    # 2. Enviar comando al xArm Esclavo
    arm_slave.???                 # <-- completar

    # 3. Leer torques articulares
    _, torques = arm_slave.???    # <-- completar

    # 4. Estimar fuerza externa (Ec. 3)
    F_ext = estimate_force(torques, arm_slave.get_joint_states())

    # 5. Enviar F_ext al Maestro
    send_to_master(F_ext)         # <-- implementar

    time.sleep(1.0 / CONTROL_HZ)
\end{lstlisting}

\subsubsection*{Pregunta~C2 \pts{8}}
\begin{enumerate}[label=(\alph*)]
  \item En el robot Maestro, la retroalimentación háptica se implementa
        aplicando torques articulares proporcionales a la fuerza estimada:
        \[
          \bm{\tau}_{m,fb} = \bm{J}_m^\top \cdot \beta \cdot \hat{\bm{F}}_{ext}
        \]
        ¿Qué representa físicamente $\beta$?
        ¿Cuál es el valor máximo de $\beta$ que permite su sistema
        antes de volverse inestable?
        Justifiquen con la condición de pasividad de la red háptica.
  \item ¿El xArm Lite 6 permite aplicar torques articulares directamente
        desde la SDK? Investiguen el modo \texttt{SERVO} o
        \texttt{TORQUE} de la API~\cite{xarm_github}.
        Si no es posible directamente, ¿qué alternativa implementaron
        para simular el efecto háptico en el maestro?
  \item Muestren el fragmento de código Python que implementa el
        feedback háptico en el maestro.
\end{enumerate}

\TODO{Respondan aquí la Pregunta C2. Incluyan los cálculos de $\beta_{max}$
y el código correspondiente.}

%% ============================================================
\section{Implementación con xArm SDK}
\label{sec:implementation}
%% ============================================================

\begin{theorybox}[xArm Python SDK]
  El repositorio oficial~\cite{xarm_github} provee la librería
  \texttt{xArm-Python-SDK} con la clase \texttt{XArmAPI}.
  Los modos de operación relevantes son:
  \textbf{Modo~0} (posición), \textbf{Modo~1} (servo/ángulo),
  \textbf{Modo~4} (velocidad articular), \textbf{Modo~6} (pose cartesiana).
  El estado del robot se controla con \texttt{set\_state()}:
  estado~0 (sport), estado~3 (pausa), estado~4 (stop).
\end{theorybox}

\subsubsection*{Pregunta~I1 \pts{10}}
\begin{enumerate}[label=(\alph*)]
  \item Documenten el procedimiento completo de configuración de
        \textbf{ambos xArm Lite 6} para su experimento:
        dirección IP de cada robot, modo de operación, estado inicial,
        parámetros de seguridad (velocidad máxima, torque máximo,
        zona de colisión).
  \item Describan el procedimiento de \textbf{calibración} que realizaron:
        ¿cómo determinaron la posición inicial del maestro y el esclavo?
        ¿Cómo sincronizaron sus espacios de trabajo?
  \item Encontraron algún \textbf{bug o limitación} de la SDK que
        debieron superar? Documenten la solución implementada.
        (Ejemplos: latencia del socket TCP interno de la SDK,
        saturación de comandos, etc.)
  \item ¿Cómo manejaron las \textbf{excepciones de seguridad}
        (\textit{error codes}) del xArm? Escriban el código de
        manejo de errores implementado.
\end{enumerate}

\TODO{Respondan aquí la Pregunta I1. Esta sección debe ser suficientemente
detallada para que otro equipo pueda replicar exactamente su setup.}

\begin{lstlisting}[caption={Configuración inicial de los xArm — Completar con su implementación},
                   label={lst:setup}]
# ── Completar con el código real de configuración ──
from xarm.wrapper import XArmAPI

def setup_robot(ip, max_vel=None, max_acc=None):
    """Configurar xArm para teleoperación.
    
    Args:
        ip: Dirección IP del robot (str)
        max_vel: Velocidad máxima articular (°/s)
        max_acc: Aceleración máxima articular (°/s²)
    Returns:
        arm: instancia configurada de XArmAPI
    """
    arm = XArmAPI(ip)
    arm.motion_enable(enable=True)
    arm.set_mode(?)        # <-- justificar
    arm.set_state(0)
    
    # Configuración de límites de seguridad
    arm.set_joint_maxacc(?)    # <-- valor elegido
    arm.set_tcp_maxacc(?)      # <-- valor elegido
    
    # Posición inicial (home)
    arm.set_servo_angle(angle=?, speed=?, wait=True)
    
    return arm

# Inicializar ambos robots
arm_master = setup_robot('192.168.1.???')  # <-- IP real maestro
arm_slave  = setup_robot('192.168.1.???')  # <-- IP real esclavo
\end{lstlisting}

%% ============================================================
\section{Resultados Experimentales}
\label{sec:experiments}
%% ============================================================

\subsection{Configuración Experimental}

\subsubsection*{Pregunta~E1 \pts{6}}
Describan con precisión el setup experimental:
\begin{enumerate}[label=(\alph*)]
  \item Especificaciones del hardware utilizado:
  \smallskip
  \resizebox{\linewidth}{!}{%
  \begin{tabular}{lll}
    \toprule
    \textbf{Componente} & \textbf{PC Maestro} & \textbf{PC Esclavo} \\
    \midrule
    CPU      & \answerline[2cm] & \answerline[2cm] \\
    RAM      & \answerline[2cm] & \answerline[2cm] \\
    OS       & \answerline[2cm] & \answerline[2cm] \\
    Python   & \answerline[2cm] & \answerline[2cm] \\
    xArm SDK & \answerline[2cm] & \answerline[2cm] \\
    \bottomrule
  \end{tabular}}
  \smallskip
  \item Describan la \textbf{tarea experimental} realizada
        (trayectoria seguida, objeto con el que interactuó el esclavo,
        tipo de contacto generado deliberadamente).
  \item ¿Cuántos ensayos realizaron? ¿Cuál fue el criterio de éxito
        de cada ensayo?
\end{enumerate}

\TODO{Respondan aquí la Pregunta E1.}

\subsection{Desempeño de Seguimiento de Posición}

\subsubsection*{Pregunta~E2 \pts{12}}
\begin{enumerate}[label=(\alph*)]
  \item Grafiquen el \textbf{error de seguimiento articular}
        $e_i(t) = q_i^{master}(t) - q_i^{slave}(t)$ para
        las 6 articulaciones durante un ensayo representativo.
        ¿Cuál articulación presentó mayor error? ¿Por qué?
  \item Calculen las métricas de desempeño y llenen la tabla:

  \smallskip
  \resizebox{\linewidth}{!}{%
  \begin{tabular}{lccc}
    \toprule
    \textbf{Articulación} & \textbf{Error RMS (°)} & \textbf{Error máx. (°)} & \textbf{Retardo (ms)} \\
    \midrule
    Joint 1 & \answerline[1.5cm] & \answerline[1.5cm] & \answerline[1.5cm] \\
    Joint 2 & \answerline[1.5cm] & \answerline[1.5cm] & \answerline[1.5cm] \\
    Joint 3 & \answerline[1.5cm] & \answerline[1.5cm] & \answerline[1.5cm] \\
    Joint 4 & \answerline[1.5cm] & \answerline[1.5cm] & \answerline[1.5cm] \\
    Joint 5 & \answerline[1.5cm] & \answerline[1.5cm] & \answerline[1.5cm] \\
    Joint 6 & \answerline[1.5cm] & \answerline[1.5cm] & \answerline[1.5cm] \\
    \midrule
    \textbf{Promedio} & \answerline[1.5cm] & \answerline[1.5cm] & \answerline[1.5cm] \\
    \bottomrule
  \end{tabular}}
  \smallskip

  \item El retardo de seguimiento medido, ¿es consistente con el RTT
        de red medido en la Pregunta A2? Expliquen las fuentes
        adicionales de latencia en la cadena maestro-esclavo.
  \item ¿En qué momentos el error fue mayor? ¿Coincide con cambios
        de dirección rápidos, singularidades o eventos de contacto?
\end{enumerate}

\FIGURE{Fig. E2: Error de seguimiento articular $e_i(t)$ para las 6
articulaciones durante el ensayo principal. Señalen los instantes
de contacto.}

\TODO{Respondan aquí la Pregunta E2.}

\subsection{Estimación y Retroalimentación de Fuerza}

\subsubsection*{Pregunta~E3 \pts{14} — \textit{Sección central del entregable}}
\begin{enumerate}[label=(\alph*)]
  \item Grafiquen $\hat{\bm{F}}_{ext}(t) = [F_x, F_y, F_z]^\top$
        durante un evento de contacto controlado
        (e.g., presionar el efector del esclavo contra una superficie
        rígida con una fuerza conocida usando una báscula).
        ¿Cuál fue la fuerza de contacto máxima registrada?
  \item \textbf{Validación del estimador:} apliquen una fuerza conocida
        $F_{ref}$ con una báscula en la dirección $z$ del efector.
        Grafiquen $\hat{F}_z(t)$ vs. $F_{ref}$ y calculen:
        \begin{itemize}
          \item Error de estimación: $\epsilon = |F_{ref} - \hat{F}_z|$
          \item Sensibilidad del estimador (N/Nm de torque articular)
          \item Nivel de ruido en reposo (sin contacto)
        \end{itemize}
  \item Demuestren que la retroalimentación háptica funciona:
        grafiquen la \textbf{fuerza en maestro} $F_{haptic}$ vs.
        la \textbf{fuerza estimada en esclavo} $\hat{F}_{slave}$.
        ¿Son proporcionales con factor $\beta$?
  \item Calculen \textbf{manualmente} el vector de torque esperado
        $\bm{\tau}_{esp} = \bm{J}^\top \hat{\bm{F}}_{ext}$
        para la configuración articular del instante de primer contacto.
        Comparen con $\bm{\tau}_{meas}$ de sus gráficas.
        ¿Coinciden en orden de magnitud?
\end{enumerate}

\FIGURE{Fig. E3a: Fuerza estimada $\hat{\bm{F}}_{ext}(t) = [F_x, F_y, F_z]^\top$
durante el evento de contacto. Señalen: inicio de contacto ($F > F_{th}$),
fuerza máxima, y fin de contacto.}

\FIGURE{Fig. E3b: Validación del estimador: fuerza de referencia
$F_{ref}$ (báscula) vs. fuerza estimada $\hat{F}_z$ en función del tiempo.
Incluyan la recta de regresión lineal y $R^2$.}

\TODO{Respondan aquí la Pregunta E3. Esta es la sección más importante
del entregable. Muestren evidencia fotográfica del setup de validación
con la báscula.}

\subsection{Análisis de Estabilidad del Sistema Háptico}

\subsubsection*{Pregunta~E4 \pts{8}}
\begin{enumerate}[label=(\alph*)]
  \item Realizaron la prueba de \textbf{contacto con pared rígida}
        (hard wall contact) tanto en simulación (pared simulada con
        alto $K_d$) como en objeto real. ¿El sistema fue estable en
        ambos casos? Grafiquen la respuesta de fuerza.
  \item Identifiquen el valor de $\beta_{crit}$ experimental
        (incrementando $\beta$ hasta observar oscilaciones).
        ¿Coincide con el valor teórico calculado en C2(a)?
  \item ¿Observaron algún \textbf{acoplamiento de inestabilidad}
        cuando el maestro y el esclavo estaban en contacto simultáneo
        con sus respectivas superficies? Describan el fenómeno.
\end{enumerate}

\TODO{Respondan aquí la Pregunta E4.}

%% ============================================================
\section{Discusión}
\label{sec:discussion}
%% ============================================================

\subsubsection*{Pregunta~D1 \pts{8}}
\begin{enumerate}[label=(\alph*)]
  \item \textbf{Limitaciones del sistema implementado:}
        Identifiquen y analicen cuantitativamente al menos \textbf{tres
        limitaciones} de su implementación (e.g., precisión del
        estimador de fuerza, latencia, espacio de trabajo, etc.).
        Para cada limitación, propongan una solución concreta y
        viable con el hardware disponible.
  \item \textbf{Comparación con estado del arte:}
        Comparen sus métricas de desempeño (error de posición,
        precisión de fuerza, latencia) con al menos DOS trabajos
        publicados de teleoperación con brazos robóticos similares.
        ¿En qué aspectos su sistema es competitivo? ¿En cuáles
        queda rezagado?
  \item \textbf{Escalabilidad:} Si tuvieran que operar este sistema
        con un retardo de $T_d = 200$~ms (e.g., operación
        transcontinental), ¿qué modificaciones al controlador
        serían necesarias? Mencionen al menos una técnica
        del estado del arte (e.g., wave variables, model mediation).
  \item \textbf{Aplicación real:} Propongan una aplicación concreta
        en la industria o medicina donde su sistema, con las mejoras
        indicadas, podría tener impacto. Estimen el orden de magnitud
        del costo de implementación.
\end{enumerate}

\TODO{Respondan aquí la Pregunta D1. Esta sección demuestra la
profundidad de comprensión del equipo. Usen referencias bibliográficas.}

%% ============================================================
\section{Conclusiones}
\label{sec:conclusions}
%% ============================================================

\subsubsection*{Pregunta~Cn1 \pts{6}}
Redacten una sección de conclusiones que responda:
\begin{enumerate}[label=(\alph*)]
  \item ¿Lograron implementar teleoperación con sensado de fuerza
        funcional entre los dos xArm Lite 6? ¿Con qué nivel de
        desempeño (cuantifiquen)?
  \item ¿Cuál fue el mayor desafío técnico encontrado y cómo lo resolvieron?
  \item ¿Qué aprendizaje clave obtuvieron sobre la relación entre
        el Jacobiano, los torques articulares y la fuerza cartesiana?
  \item ¿Qué harían diferente si repitieran el proyecto?
\end{enumerate}

\TODO{Escriban aquí las conclusiones del equipo en forma de párrafos
continuos (no lista). Extensión: 200--300 palabras.}

%% ============================================================
\section*{Agradecimientos}
%% ============================================================

\TODO{Agradezcan al Computational Robotics Lab, al Prof. Alberto Muñoz,
y a cualquier persona o institución que haya apoyado su trabajo.}

%% ============================================================
%% BIBLIOGRAFÍA
%% ============================================================
\begin{thebibliography}{10}

\bibitem{xarm_github}
  UFACTORY, ``xArm-Python-SDK,''
  GitHub repository, 2024.
  [Online]. Available: \url{https://github.com/xArm-developer/xArm-Python-SDK}

\bibitem{siciliano2009}
  B. Siciliano, L. Sciavicco, L. Villani, G. Oriolo,
  \textit{Robotics: Modelling, Planning and Control},
  Springer, 2009.

\bibitem{spong2006}
  M.W. Spong, S. Hutchinson, M. Vidyasagar,
  \textit{Robot Modeling and Control},
  Wiley, 2006.

\bibitem{hogan1985}
  N. Hogan,
  ``Impedance Control: An Approach to Manipulation,''
  \textit{ASME J. Dyn. Sys., Meas., Control},
  vol. 107, 1985.

\bibitem{colgate1993}
  J.E. Colgate, J.M. Brown,
  ``Factors affecting the Z-width of a haptic display,''
  in \textit{Proc. IEEE ICRA}, 1994, pp. 3205--3210.

\bibitem{murray1994}
  R.M. Murray, Z. Li, S.S. Sastry,
  \textit{A Mathematical Introduction to Robotic Manipulation},
  CRC Press, 1994.

\bibitem{niemeyer2008}
  G. Niemeyer, C. Preusche, G. Hirzinger,
  ``Telerobotics,''
  in \textit{Springer Handbook of Robotics},
  Springer, 2008.

\bibitem{thompson2001}
  S. Thompson, I. Reid, \textbf{A. Muñoz}, D.W. Murray,
  ``Using a stationary camera to improve the performance
  of a teleoperated robot,''
  \textit{IEEE Trans. Systems, Man \& Cybernetics-A},
  vol. 31(1), pp. 74--82, 2001.

\bibitem{ref9}
  %% COMPLETAR: Citar un artículo IEEE de los últimos 5 años sobre
  %% teleoperación con sensado de fuerza.
  %% Buscar en IEEE Xplore: "force-sensing teleoperation robot arm 2020-2025"
  A.~Autor, B.~Autor, ``Título del artículo,''
  \textit{IEEE Trans. Robotics}, vol.~XX, no.~Y, pp.~ZZ--ZZ, 202X.

\bibitem{ref10}
  %% COMPLETAR: Citar un artículo IEEE reciente sobre estimación
  %% de fuerza sin sensor externo (torque-based force estimation).
  C.~Autor, D.~Autor, ``Título del artículo,''
  in \textit{Proc. IEEE ICRA}, 202X, pp.~ZZ--ZZ.

\end{thebibliography}

%% ============================================================
%% APÉNDICE — Rúbrica de Evaluación Completa
%% ============================================================
\appendix
\section{Rúbrica de Evaluación}
\label{app:rubrica}

\begin{center}
\small
\begin{tabular}{p{0.35\columnwidth}ccp{0.25\columnwidth}}
  \toprule
  \textbf{Pregunta / Criterio} & \textbf{Pts.} & \textbf{Mín.} & \textbf{Evidencia Requerida} \\
  \midrule
  \multicolumn{4}{l}{\textbf{Teoría y Fundamentos}} \\
  T1: Jacobiano y singularidades xArm & 8 & 5 & Derivación + figura sing. \\
  T2: Control de par calculado (CTC)  & 10 & 6 & Derivación + gráfica resp. escalón \\
  T3: Control de impedancia           & 8 & 5 & Cálculo $K_{d,max}$ + tabla $K_d$ \\
  \midrule
  \multicolumn{4}{l}{\textbf{Arquitectura y Comunicación}} \\
  A1: Diagrama de bloques completo    & 10 & 6 & Figura con señales dimensionadas \\
  A2: Protocolo y métricas de red     & 8 & 5 & Tabla RTT/jitter + captura ping \\
  \midrule
  \multicolumn{4}{l}{\textbf{Sensado de Fuerza}} \\
  F1: Estimador torque-fuerza         & 12 & 7 & Gráficas antes/después filtro \\
  F2: Comparativa estrategias         & 8 & 5 & Tabla completada + justificación \\
  \midrule
  \multicolumn{4}{l}{\textbf{Control e Implementación}} \\
  C1: Control del esclavo             & 10 & 6 & Código funcional + flowchart \\
  C2: Feedback háptico en maestro     & 8 & 5 & Cálculo $\beta_{max}$ + código \\
  I1: Setup y configuración xArm SDK  & 10 & 6 & Documentación reproducible \\
  \midrule
  \multicolumn{4}{l}{\textbf{Resultados Experimentales}} \\
  E1: Setup y descripción experimentos & 6 & 4 & Tabla hardware + descripción \\
  E2: Error de seguimiento (6 joints)  & 12 & 7 & Gráficas + tabla métricas \\
  E3: Estimación y validación fuerza   & 14 & 9 & Gráficas + validación báscula \\
  E4: Análisis estabilidad háptica     & 8 & 5 & Prueba hard-wall + $\beta_{crit}$ \\
  \midrule
  \multicolumn{4}{l}{\textbf{Análisis y Conclusiones}} \\
  D1: Discusión limitaciones y EA     & 8 & 5 & Citas IEEE + comparativa cuant. \\
  Cn1: Conclusiones                   & 6 & 4 & Párrafos 200--300 palabras \\
  \midrule
  \multicolumn{4}{l}{\textbf{Calidad del Documento IEEE}} \\
  Formato y ortografía                & 4 & 2 & LaTeX compilable, sin errores \\
  Figuras con pie de figura y etiquetas & 4 & 2 & Mín. 8 figuras propias \\
  Referencias (mín. 10, 2 recientes)  & 4 & 2 & BibTeX correcto \\
  \midrule
  \rowcolor{tec_blue!15}
  \textbf{TOTAL} & \textbf{160} & \textbf{96} & \\
  \bottomrule
\end{tabular}
\end{center}

\begin{rubricbox}[Criterio de aprobación]
  \begin{itemize}[itemsep=0.2em]
    \item Score $\geq 96$ puntos sobre 160 (60\%).
    \item La sección E3 debe incluir \textbf{evidencia experimental real}
          con los xArm físicos (no simulación).
    \item El estimador de fuerza debe mostrar coherencia con los torques
          articulares medidos: $\bm{\tau}_{meas} \approx \bm{J}^\top \hat{\bm{F}}_{ext}$.
    \item El documento LaTeX debe compilar sin errores.
    \item Mínimo \textbf{8 figuras originales} (no tomadas de internet).
  \end{itemize}
\end{rubricbox}

\section{Guía de Entrega}
\label{app:entrega}

\begin{todobox}[Checklist de entrega]
  Antes de entregar, verifiquen que el repositorio/zip contiene:
  \begin{enumerate}[label=$\square$, itemsep=0.3em]
    \item \texttt{reporte.pdf} — artículo IEEE compilado (este documento completo)
    \item \texttt{reporte.tex} — fuente LaTeX del artículo
    \item \texttt{master\_control.py} — código del robot Maestro
    \item \texttt{slave\_control.py} — código del robot Esclavo
    \item \texttt{force\_estimator.py} — módulo de estimación de fuerza
    \item \texttt{net\_test.py} — herramienta de test de red (adaptado para xArm)
    \item \texttt{data/} — carpeta con los archivos \texttt{.csv} o \texttt{.npy}
          de cada ensayo experimental (mínimo 3 ensayos)
    \item \texttt{figures/} — imágenes originales (PNG/PDF) usadas en el reporte
    \item \texttt{video\_demo.mp4} — video de 60--120 s mostrando el sistema
          funcionando en tiempo real con ambos xArm
    \item \texttt{README.md} — instrucciones para reproducir los experimentos
  \end{enumerate}
\end{todobox}

\begin{rubricbox}[Video de demostración — Puntos extra]
  Un video de demostración de alta calidad que muestre:
  \begin{itemize}[itemsep=0.2em]
    \item El operador moviendo el Maestro y el Esclavo replicando el movimiento.
    \item Un evento de contacto claro con las gráficas de fuerza en pantalla.
    \item La reacción háptica perceptible en el Maestro.
  \end{itemize}
  Un video bien editado (subtítulos, señales anotadas) puede obtener
  hasta \textbf{+10 puntos extra} a criterio del profesor.
\end{rubricbox}

\end{document}
%% ============================================================
%% FIN DEL DOCUMENTO
%% ============================================================
